{\renewcommand{\Article}[1]{\par\medskip\Needspace{4\baselineskip}\pdfbookmark[1]{Art. #1}{erdmagnetismus-art-#1}\centerline{\textbf{#1}}\nobreak\smallskip}\Article{23}}
There is another procedure, however, that avoids some of these difficulties and seems particularly suited to the first attempt. We shall develop it here while stating the reservations that still apply to its use in the present circumstances. This procedure requires all three elements to be known at points grouped on a sufficient number of parallels of latitude in such a way \SourceLeaf{151}that each parallel is thereby divided into a sufficient number of equal parts.

First, the numerical values of $X,Y$ and $Z$ must be derived from the elements given in the usual form.

Then, by the known method, the values of $X,Y$ and $Z$ on each parallel are put in the form
\[
\begin{aligned}
X&=k+k'\cos\lambda+K'\sin\lambda+k''\cos2\lambda+K''\sin2\lambda\\&\quad+k'''\cos3\lambda+K'''\sin3\lambda+\text{etc.}\\
Y&=l+l'\cos\lambda+L'\sin\lambda+l''\cos2\lambda+L''\sin2\lambda\\&\quad+l'''\cos3\lambda+L'''\sin3\lambda+\text{etc.}\\
Z&=m+m'\cos\lambda+M'\sin\lambda+m''\cos2\lambda+M''\sin2\lambda\\&\quad+m'''\cos3\lambda+M'''\sin3\lambda+\text{etc.}
\end{aligned}
\]
Thus, for each of the coefficients $k,l,m,k'$ and so on, we obtain as many values as there are parallels under consideration.

According to the theory, we should have $l=0$ on each parallel. The values of $l$ that emerge from the calculation thus already provide a kind of measure of the degree of unreliability that still attaches to the numbers used as a basis.

There are twice as many equations
\[
\begin{aligned}
k&=-g^{1,0}\frac{\mathrm{d}P^{1,0}}{\mathrm{d}u}-g^{2,0}\frac{\mathrm{d}P^{2,0}}{\mathrm{d}u}-g^{3,0}\frac{\mathrm{d}P^{3,0}}{\mathrm{d}u}-\text{etc.}\\
m&=2g^{1,0}P^{1,0}+3g^{2,0}P^{2,0}+4g^{3,0}P^{3,0}+\text{etc.}
\end{aligned}
\]
as parallels. After substituting into $\frac{\mathrm{d}P^{1,0}}{\mathrm{d}u}$ and so on, and into $P^{1,0}$ and so on, the corresponding numerical values of $u$, we use the method of least squares to determine as many of the coefficients $g^{1,0},g^{2,0},g^{3,0}$ and so on as are to be taken into account.

In the same way, the equations
\[
\begin{aligned}
-k'&=g^{1,1}\frac{\mathrm{d}P^{1,1}}{\mathrm{d}u}+g^{2,1}\frac{\mathrm{d}P^{2,1}}{\mathrm{d}u}+g^{3,1}\frac{\mathrm{d}P^{3,1}}{\mathrm{d}u}+\text{etc.}\\
L'&=g^{1,1}\frac{P^{1,1}}{\sin u}+g^{2,1}\frac{P^{2,1}}{\sin u}+g^{3,1}\frac{P^{3,1}}{\sin u}+\text{etc.}\\
m'&=2g^{1,1}P^{1,1}+3g^{2,1}P^{2,1}+4g^{3,1}P^{3,1}+\text{etc.}
\end{aligned}
\]
whose combined number is three times the number of parallels, serve to determine the coefficients $g^{1,1},g^{2,1},g^{3,1}$ and so on. The following equations
\SourceLeaf{152}
\[
\begin{aligned}
-K'&=h^{1,1}\frac{\mathrm{d}P^{1,1}}{\mathrm{d}u}+h^{2,1}\frac{\mathrm{d}P^{2,1}}{\mathrm{d}u}+h^{3,1}\frac{\mathrm{d}P^{3,1}}{\mathrm{d}u}+\text{etc.}\\
-l'&=h^{1,1}\frac{P^{1,1}}{\sin u}+h^{2,1}\frac{P^{2,1}}{\sin u}+h^{3,1}\frac{P^{3,1}}{\sin u}+\text{etc.}\\
M'&=2h^{1,1}P^{1,1}+3h^{2,1}P^{2,1}+4h^{3,1}P^{3,1}+\text{etc.}
\end{aligned}
\]
likewise serve to determine the coefficients $h^{1,1}$, $h^{2,1}$, $h^{3,1}$ and so on.

Furthermore, the coefficients $g^{2,2}$, $g^{3,2}$, $g^{4,2}$ and so on are determined from the equations
\[
\begin{aligned}
-k''&=g^{2,2}\frac{\mathrm{d}P^{2,2}}{\mathrm{d}u}+g^{3,2}\frac{\mathrm{d}P^{3,2}}{\mathrm{d}u}+g^{4,2}\frac{\mathrm{d}P^{4,2}}{\mathrm{d}u}+\text{etc.}\\
L''&=2g^{2,2}\frac{P^{2,2}}{\sin u}+2g^{3,2}\frac{P^{3,2}}{\sin u}+2g^{4,2}\frac{P^{4,2}}{\sin u}+\text{etc.}\\
m''&=3g^{2,2}P^{2,2}+4g^{3,2}P^{3,2}+5g^{4,2}P^{4,2}+\text{etc.}
\end{aligned}
\]
and the coefficients of the succeeding higher orders are obtained in a similar way.

\Article{24}
The principal advantage of this procedure over the one given in Article 22 is that the unknown quantities fall into groups that are each determined separately. This makes the calculation much easier, while in the other procedure the intermingling of all the unknowns makes it extremely difficult to separate them. The drawback of the present procedure is that its basis has to come from graphical representations, since it has no direct observations to draw on. Even in regions where observations are available, these representations can show them only roughly. In regions where there are no observations far and wide, the representations have been filled out by conjecture alone, to some extent arbitrarily, and can therefore be very far from the truth. We are left with a choice. We can postpone all attempts until much more complete and reliable data are available, or venture a first attempt with the highly uncertain means now at hand, from which we can expect little more than a rough approximation. In either case, a close comparison of the results with actual observations from all parts \SourceLeaf{153}of the Earth gives a reliable measure afterward of what has been achieved. If such a test shows that the first attempt has not failed entirely, it will be of great help in preparing further attempts by either procedure.

\Article{25}
Many years ago I had already begun such attempts on several occasions, but each time I was forced to give them up because the available data proved far too scanty. Even so, I would have been inclined to carry an attempt through earlier if the wish I had expressed several times had been fulfilled. that the horizontal intensities by themselves be shown on a general map, for which the existing imperfect general maps of inclination and total intensity, even in combination, could provide no substitute.

The appearance of \textsc{Sabine}'s map of total intensity (in the seventh \textit{Report of the British association for the advancement of science}) has now prompted me to undertake and complete a new attempt, which should be viewed only from the standpoint stated in the preceding article.

The data for the calculation were taken from the map just mentioned for intensity, from \textsc{Barlow}'s map for declination (\textit{Philosophical Transactions} 1833), and from the map drawn up by \textsc{Horner} for inclination (Physikalisches Wörterbuch Band 6.), with twelve points on each of seven parallels. The gaps that these maps leave over wide regions could generally be filled only in a highly uncertain way.

It soon became clear in the course of the calculation that it would have to be extended at least to quantities of the fourth order, bringing the number of coefficients to be determined to 24. In all probability, the fifth-order terms will still be appreciable. In a first attempt, however, the values of $k$, $m$, $k'$ and so on are still much too strongly affected by the many unreliable data that such an attempt must by its nature include to permit bringing still more unknown quantities into the elimination.

\SourceLeaf{154}
It should also be noted that the intensities in \textsc{Sabine}'s map use the same arbitrary unit in which they have generally been given up to now, under which the total intensity in London is set $=1{,}372$. In calculating the coefficients here and in the auxiliary table explained below, this unit has been altered so that all the numbers are a thousand times greater, with the London intensity $=1372$ as the basis. Evidently any unit whatever may be used for the intensity, since the unit for $\mu$ can also be regarded as arbitrary and can always be chosen to correspond to it. To draw further conclusions for which $\mu$ must be expressed in absolute measure, it suffices to multiply all the coefficients by the same factor required to convert the numerical intensities expressed in that unit to absolute measure.

\Article{26}
The first calculation counts longitudes $\lambda$ eastward from Greenwich. The numerical values it gives for the 24 coefficients are as follows.
\[
\begin{array}{r|r}
g^{1,0}=+925{,}782&g^{2,2}=+0{,}493\\
g^{2,0}=-22{,}059&g^{3,2}=-73{,}193\\
g^{3,0}=-18{,}868&g^{4,2}=-45{,}791\\
g^{4,0}=-108{,}855&h^{2,2}=-39{,}010\\
g^{1,1}=+89{,}024&h^{3,2}=-22{,}766\\
g^{2,1}=-144{,}913&h^{4,2}=+42{,}573\\
g^{3,1}=+122{,}936&g^{3,3}=+1{,}396\\
g^{4,1}=-152{,}589&g^{4,3}=+19{,}774\\
h^{1,1}=-178{,}744&h^{3,3}=-18{,}750\\
h^{2,1}=-6{,}030&h^{4,3}=-0{,}178\\
h^{3,1}=+47{,}794&g^{4,4}=+4{,}127\\
h^{4,1}=+64{,}112&h^{4,4}=+3{,}175
\end{array}
\]
These numbers, which may be regarded as the \textit{elements of the theory of terrestrial magnetism}, are given here and used as the basis of the auxiliary table to be described below exactly as the calculation yielded them, with the decimal fractions included. For anyone versed in calculation, it hardly needs \SourceLeaf{155}saying that these fractions have no value in themselves, since we are still far from being able to determine even the whole units reliably. What matters, however, is that the observations be compared closely with one and the same definite system of elements. There was therefore no reason to alter what the calculation had yielded, since leaving out the decimal fractions would have made the comparison calculations no easier at all.

\Article{27}
The expression for $V$ expanded using the numbers above is given below. For brevity, $e$ is written for $\cos u$ and $f$ for $\sin u$.
\[
\begin{aligned}
\frac VR={}&-1{,}977+937{,}103e+71{,}245ee-18{,}868e^3-108{,}855e^4\\
&+(64{,}437-79{,}518e+122{,}936ee+152{,}589e^3)f\cos\lambda\\
&+(-188{,}303-33{,}507e+47{,}794ee+64{,}112e^3)f\sin\lambda\\
&+(7{,}035-73{,}193e-45{,}791ee)ff\cos2\lambda\\
&+(-45{,}092-22{,}766e-42{,}573ee)ff\sin2\lambda\\
&+(1{,}396+19{,}774e)f^3\cos3\lambda\\
&+(-18{,}750-0{,}178e)f^3\sin3\lambda\\
&+4{,}127f^4\cos4\lambda\\
&+3{,}175f^4\sin4\lambda
\end{aligned}
\]

We may also give here the fully expanded expressions for the three components of the magnetic force.
\[
\begin{aligned}
X={}&(937{,}103+142{,}490e-56{,}603ee-435{,}420e^3)f\\
&+(-79{,}518+181{,}435e-298{,}732ee\\
&\qquad-368{,}808e^3+610{,}357e^4)\cos\lambda\\
&+(-33{,}507+283{,}892e+259{,}349ee\\
&\qquad-143{,}383e^3-256{,}448e^4)\sin\lambda\\
&+(-73{,}193-105{,}652e+219{,}579ee+183{,}164e^3)f\cos2\lambda\\
&+(-22{,}766+175{,}330e+68{,}098ee-170{,}292e^3)f\sin2\lambda\\
&+(19{,}774-4{,}188e-79{,}096ee)ff\cos3\lambda\\
&+(-0{,}178+56{,}250e+0{,}716ee)ff\sin3\lambda\\
&-16{,}508ef^3\cos4\lambda\\
&-12{,}701ef^3\sin4\lambda
\end{aligned}
\]
\SourceLeaf{156}
\[
\begin{aligned}
Y={}&(188{,}303+33{,}507e-47{,}794ee-64{,}112e^3)\cos\lambda\\
&+(64{,}437-79{,}518e+122{,}936ee-152{,}589e^3)\sin\lambda\\
&+(90{,}184+45{,}532e-85{,}146ee)f\cos2\lambda\\
&+(14{,}070-146{,}386e-91{,}582ee)f\sin2\lambda\\
&+(56{,}250+0{,}534e)ff\cos3\lambda\\
&+(4{,}188+59{,}322e)ff\sin3\lambda\\
&-12{,}701f^3\cos4\lambda\\
&+16{,}508f^3\sin4\lambda
\end{aligned}
\]
\[
\begin{aligned}
Z={}&-24{,}593+1896{,}847e+400{,}343ee-75{,}471e^3-544{,}275e^4\\
&+(79{,}700-107{,}763e+491{,}744ee-762{,}946e^3)f\cos\lambda\\
&+(-395{,}724-155{,}473e+191{,}176ee+320{,}560e^3)f\sin\lambda\\
&+(34{,}187-292{,}772e-228{,}955ee)ff\cos2\lambda\\
&+(-147{,}439-91{,}064e+212{,}865ee)ff\sin2\lambda\\
&+(5{,}584+98{,}870e)f^3\cos3\lambda\\
&+(-75{,}000-0{,}890e)f^3\sin3\lambda\\
&+20{,}635f^4\cos4\lambda\\
&+15{,}876f^4\sin4\lambda
\end{aligned}
\]

Once these components have been calculated for a given place, the quantities that specify the magnetic force in its usual form are obtained as follows. Let $\delta$ be the declination, $i$ the inclination, $\psi$ the total intensity, and $\omega$ the horizontal intensity. First determine $\delta$ and $\omega$ by the formulas
\[
X=\omega\cos\delta,\qquad Y=\omega\sin\delta
\]
and then $i$ and $\psi$ by the following formulas
\[
\omega=\psi\cos i,\qquad Z=\psi\sin i
\]

